%=====================================================================
%   LITERATURE REVIEW
%=====================================================================
%   Replace this chapter with your own. What is here is a demonstration of
%   how to cite, kept because it is the thing people get wrong.
%
%   The keys cited below are PLACEHOLDERS from thesis.bib. They exist to show
%   the citation commands working. Do not leave them in your thesis.
%=====================================================================
\chapter{Literature review}
\label{ch:review}

%---------------------------------------------------------------------
%   HOW TO CITE
%---------------------------------------------------------------------
%   natbib gives you two forms, and they are NOT interchangeable.
%
%     \citet{key}   TEXTUAL. The author becomes part of the sentence.
%                   "As \citet{shannon1948} showed, ..."
%                   -> As Shannon (1948) showed, ...
%
%     \citep{key}   PARENTHETICAL. The whole citation sits in brackets at
%                   the end of a clause.
%                   "...is a solved problem \citep{shannon1948}."
%                   -> ...is a solved problem (Shannon, 1948).
%
%   \cite{key} behaves like \citet in author-year mode and like a numeric
%   citation in IEEE mode, so it reads differently depending on the switch at
%   the top of 0_main.tex. Prefer \citet and \citep: they always mean what
%   they say.
%
%   In IEEE mode both come out as [1], which is why IEEE citations are easier
%   to write and harder to read.
%---------------------------------------------------------------------

Textual, where the author is the subject of the sentence:
\citet{shannon1948} established the information-theoretic limits that still
bound modern communication systems.

Parenthetical, at the end of a clause:
The effect is well documented \citep{placeholder_article}.

Several sources at once, which natbib sorts and compresses:
The same conclusion is reached by three independent routes
\citep{placeholder_article,placeholder_book,placeholder_conference}.

A page, chapter or section within a source:
The register map is given in full \citep[Ch.~8]{placeholder_website}.

The year alone, where the sentence already names the author:
Shannon's channel capacity theorem \citeyearpar{shannon1948} is the origin of
the result.

%---------------------------------------------------------------------
%   thesis.bib ships with one example of each entry type: article, book,
%   chapter in an edited book, conference paper, thesis, technical report and
%   website. Copy the nearest one and edit it.
%---------------------------------------------------------------------
